% Options for packages loaded elsewhere
\PassOptionsToPackage{unicode}{hyperref}
\PassOptionsToPackage{hyphens}{url}
\PassOptionsToPackage{dvipsnames,svgnames,x11names}{xcolor}
%
\documentclass[
  10pt,
]{article}
\usepackage{amsmath,amssymb}
\usepackage{iftex}
\ifPDFTeX
  \usepackage[T1]{fontenc}
  \usepackage[utf8]{inputenc}
  \usepackage{textcomp} % provide euro and other symbols
\else % if luatex or xetex
  \usepackage{unicode-math} % this also loads fontspec
  \defaultfontfeatures{Scale=MatchLowercase}
  \defaultfontfeatures[\rmfamily]{Ligatures=TeX,Scale=1}
\fi
\usepackage{lmodern}
\ifPDFTeX\else
  % xetex/luatex font selection
\fi
% Use upquote if available, for straight quotes in verbatim environments
\IfFileExists{upquote.sty}{\usepackage{upquote}}{}
\IfFileExists{microtype.sty}{% use microtype if available
  \usepackage[]{microtype}
  \UseMicrotypeSet[protrusion]{basicmath} % disable protrusion for tt fonts
}{}
\makeatletter
\@ifundefined{KOMAClassName}{% if non-KOMA class
  \IfFileExists{parskip.sty}{%
    \usepackage{parskip}
  }{% else
    \setlength{\parindent}{0pt}
    \setlength{\parskip}{6pt plus 2pt minus 1pt}}
}{% if KOMA class
  \KOMAoptions{parskip=half}}
\makeatother
\usepackage{xcolor}
\usepackage[margin=0.90in]{geometry}
\setlength{\emergencystretch}{3em} % prevent overfull lines
\providecommand{\tightlist}{%
  \setlength{\itemsep}{0pt}\setlength{\parskip}{0pt}}
\setcounter{secnumdepth}{-\maxdimen} % remove section numbering
\ifLuaTeX
  \usepackage{selnolig}  % disable illegal ligatures
\fi
\usepackage{bookmark}
\IfFileExists{xurl.sty}{\usepackage{xurl}}{} % add URL line breaks if available
\urlstyle{same}
\hypersetup{
  colorlinks=true,
  linkcolor={Maroon},
  filecolor={Maroon},
  citecolor={Blue},
  urlcolor={Blue},
  pdfauthor={Leo Gao},pdfcreator={LaTeX via pandoc}}

\author{}
\date{}

\begin{document}

\section{The exact maximum for square-free hypercube graphs with
quadratic direction
rules}\label{the-exact-maximum-for-square-free-hypercube-graphs-with-quadratic-direction-rules}

\textbf{Leo Gao} $\cdot$ AI-assisted research, 27 November 2026. \textbf{The unrestricted Erdős
problem \#86 remains open.}

Let \(Q_n\) have vertex set \(V=\mathbb F_2^n\) and basis directions
\(e_i\). A graph on the cube edges has direction indicators
\(f_i:V\to\{0,1\}\) satisfying \(f_i(x+e_i)=f_i(x)\). Let
\(p_i=\mathbb{E} f_i\), let \(p=n^{-1}\sum_i p_i\) denote its density in
the cube, and write \(b_i=2p_i-1\). The original question asks whether
the largest square-free density tends to \(1/2\).

\subsection{Main theorem}\label{main-theorem}

If the graph is square-free and \textbf{each direction indicator is a
polynomial of degree at most two over \(\mathbb F_2\)}, then \[
\sum_{i:\,b_i>0} b_i^2\le 1,\qquad
p \le \frac12+\frac1{2\sqrt n}.
\] There are no requirements on common factors, rank, symmetry, or where
the coefficients of different direction rules come from. In fact, if
\(L=4^{\lfloor\log_4 n\rfloor}\) and \(t=\lfloor(4L-n)/3\rfloor\), the
\textbf{exact maximum density} in this subclass is \[
p_{\rm quad}(n)=\frac12+\frac{n+t}{4n\sqrt L}.
\] In particular \(p_{\rm quad}(4^k)=1/2+1/(2^{k+1})\). The proof uses a
four-ary prefix tree to build a matching example in every dimension. No
unrestricted upper or lower record is claimed. Indeed \(n=9\) separates
the problems: the formula gives \textbf{1,504} edges, while Wang's fully
frustrated result gives \textbf{1,532} selected edges (\(2,304-772\)).
Thus in this dimension some nonquadratic rule is necessary to reach the
known construction.

\subsubsection{1. A property of biased
quadratics}\label{a-property-of-biased-quadratics}

If \(f\) is a quadratic Boolean function on a finite affine binary
space, let \(s=2f-1\). If \(\mathbb E s>0\), then on \textbf{each
nonempty affine hyperplane} the conditional mean of \(s\) is
nonnegative. This includes constant and affine functions.

To prove it, choose an origin, write \(B(h,x)=f(x+h)+f(x)+f(h)+f(0)\)
for the polar bilinear form, and \(K\) for its radical. For a linear
functional \(\ell\), write
\(\widehat s(\ell)=\mathbb E_x s(x)(-1)^{\ell(x)}\). Squaring and
writing the two sample points as \(x,x+h\) gives \[
\widehat s(\ell)^2=
\mathbb E_h{\bf1}_{h\in K}(-1)^{f(h)+f(0)+\ell(h)}.
\] Indeed, the average of \((-1)^{B(h,x)}\) over \(x\) vanishes outside
\(K\). The function \(h\mapsto f(h)+f(0)\) is linear on \(K\), so the
answer is either \(0\) or \(|K|/|V|\). Hence a nonzero constant Fourier
coefficient has magnitude at least that of every other coefficient. On
\(\ell(x)=a\), for nonzero \(\ell\), the conditional mean is
\(\widehat s(0)+(-1)^a\widehat s(\ell)\ge0\). The same proof applies on
affine slices.

\subsubsection{2. A missing square cannot occur in two above-half
directions}\label{a-missing-square-cannot-occur-in-two-above-half-directions}

Fix distinct \(i,j\) with \(b_i,b_j>0\). Define \[
D_i=\{x:f_i(x)=f_i(x+e_j)\},\qquad
D_j=\{x:f_j(x)=f_j(x+e_i)\}.
\] Each is the whole space or an affine hyperplane: the derivative of a
quadratic is affine, and a derivative identically one would force
balance by pairing translated points. On the complement of \(D_i\) the
translation \(x\mapsto x+e_j\) pairs opposite signs of \(s_i=2f_i-1\).
Therefore its mean on \(D_i\) is strictly positive. The same holds for
\(j\) on \(D_j\).

If \(W=D_i\cap D_j\) is nonempty, it is either all of \(D_i\) or a
hyperplane therein, and likewise for \(D_j\). Part 1, applied to the two
restricted quadratics, gives \(\mathbb E[s_i\mid W]\ge0\) and
\(\mathbb E[s_j\mid W]\ge0\). On \(W\) an \(i\)-edge and its opposite
agree, as do the \(j\)-edges. Thus the lack of a retained square gives
\(f_i+f_j\le1\) pointwise there (ordinary integers). The two conditional
means sum to at least one for the indicators, so in fact \(f_i+f_j=1\)
everywhere on \(W\). This rules out an entirely deleted square as well.
If \(W\) is empty neither kind of monochromatic square is possible.

Consequently, deleting the directions of density at most one half leaves
a two-coloring in which neither color contains a coordinate square, in
every slice. The missing-color conclusion uses quadraticity; it does not
follow from square-freeness alone.

\subsubsection{3. The weighted
inequality}\label{the-weighted-inequality}

For \(i,j\) as above, average \(s_i(x)s_j(x)\) over the four vertices of
an \(\{i,j\}\)-square. Since \(s_i\) is independent of \(e_i\) and
\(s_j\) of \(e_j\), this average is \[
\tfrac14\big(s_i(x)+s_i(x+e_j)\big)\big(s_j(x)+s_j(x+e_i)\big).
\] It vanishes unless \(x\in W\), in which case it equals \(-1\) by part
2. Hence \(\mathbb E s_i s_j=-\Pr(W)\le0\). For \(I=\{i:b_i>0\}\), set
\(Z=\sum_{i\in I} b_i s_i\). Jensen's inequality and \(s_i^2=1\) give \[
\Big(\sum_{i\in I}b_i^2\Big)^2=(\mathbb E Z)^2\le\mathbb E Z^2\le\sum_{i\in I}b_i^2.
\] It follows that \(\sum_{i\in I}b_i^2\le1\). Finally
\(2n(p-1/2)=\sum_i b_i\le\sqrt{n\sum_{i\in I} b_i^2}\le\sqrt n\).

\subsubsection{4. Optimizing the directional
biases}\label{optimizing-the-directional-biases}

The Fourier calculation in part 1 shows that every positive quadratic
bias equals \(2^{-d}\) for an integer \(d\ge0\): the polar form is
alternating and has even rank \(2d\), so \(|K|/|V|=2^{-2d}\). Thus
positive biases are the weights of at most \(n\) leaves of respective
depths \(d\) in a rooted four-ary tree. Indeed part 3 gives
\(\sum 4^{-d}\le1\), precisely the four-ary Kraft condition. For
completeness, sort the lengths increasingly and assign unused nodes at
each depth; a leaf at depth \(a\) precludes exactly \(4^{b-a}\) choices
at any deeper level \(b\), and the condition guarantees enough free
nodes. Each leaf of depth \(d\) contributes \(2^{-d}\) to the sum of
weights. Replacing a leaf by two children preserves its total weight, so
an upper bound for exactly \(n\) leaves bounds codes using fewer nonzero
biases as well.

Among trees with exactly \(n\) leaves there is an optimal one with no
unary internal vertex: suppressing one increases the weights of all its
descendants. Such trees have depth at most \(n-1\), so a maximizing tree
exists. Among maximizers choose one minimizing the sum of leaf depths.
If a shallow leaf has depth \(a\) and the deepest leaf depth is
\(b\ge a+2\), swap the shallow leaf with the parent of a deepest leaf.
That parent has \(q\ge2\) children, all leaves. Its contribution
normalized at its own depth is \(q/2\ge1\), versus \(1\) for the shallow
leaf. Moving it up cannot decrease weight, while the total depth
decreases by \((q-1)(b-1-a)>0\). Hence all leaves of a maximizing tree
occupy at most two adjacent levels.

Write \(k=\lfloor\log_4 n\rfloor\), \(L=4^k\). For \(n>L\), one cannot
fit the leaves entirely at depths at most \(k\). If all depths are at
least \(k+1\), the weight is at most \(n/(2\sqrt L)\). Otherwise the two
depths must be \(k,k+1\); the number \(u\) at depth \(k\) satisfies
\(3u+n\le4L\), hence \(u\le t\). The total weight is at most
\((n+t)/(2\sqrt L)\). For \(n=L\) the same bound follows immediately
from Cauchy--Schwarz, as it equals \(\sqrt n\). Conversely, there is a
code with \(t\) leaves at depth \(k\) and \(n-t\) at depth \(k+1\) by
the same capacity test. The maximum sum of positive biases allowed by
the inequality is exactly \((n+t)/(2\sqrt L)\).

\subsubsection{5. Realizing the upper
bound}\label{realizing-the-upper-bound}

The explicit \(Q_4\) system in the verification section has every
direction bias \(1/2\) and is square-free. Its faces have three selected
edges \emph{on average}, because its edge density is \(3/4\); each face
has at most three, so each has exactly three. Thus the product of its
four signs \(s_i=2f_i-1\) on every square is \(-1\). Call this property
full frustration.

To substitute graphs into directions of this four-direction system, put
a nonempty block of coordinates at each of its four leaves. The graph
for each block has signs \(s_a^{(i)}\); take the one-direction constant
\(+1\) system for a single leaf. At \(x\) let \(z_i\) be the parity of
block \(x_i\). Set \[
S_{i,a}(x)=s_i^{\rm out}(z_1,z_2,z_3,z_4)\,s_a^{(i)}(x_i).
\] Every edge rule is invariant in its own direction. On a within-block
square the outer factor is constant (it ignores its own parity), so the
square product is inherited from the inner graph. On a cross-block
square the inner factors cancel in pairs and the two parities flip, so
the product is inherited from the outer graph. Hence full frustration is
preserved, in particular square-freeness. The Boolean indicator of the
product is \(1+f_i^{\rm out}(z)+f_a^{(i)}(x_i)\) in \(\mathbb F_2\);
parity substitution is linear, so its degree is still at most two.
Finally \(s_i^{\rm out}\) ignores its own block, and all other nonempty
block parities are uniform independent bits. The bias of each resulting
direction is the product of its two constituent biases.

Iterate this operation on a full four-ary tree (all internal nodes have
four children). A leaf at depth \(d\) has bias \(2^{-d}\). Complete the
optimal code from part 4 by including every child at level \(k+1\) of
each non-leaf at level \(k\). There are \(4L-3t\) total leaves; this
exceeds \(n\) by at most two. Designate the unused leaves as dummy
directions and fix their bits. Averaged over all dummy assignments, the
sum of biases of the retained directions remains \((n+t)/(2\sqrt L)\).
At least one fixed assignment attains this value or better; restriction
preserves both degree at most two and square-freeness, so the upper
bound forces equality. This gives the claimed finite formula. In fact
every such assignment attains the same total by the upper bound and the
exact average.

\subsection{A different restricted theorem: a small translation
quotient}\label{a-different-restricted-theorem-a-small-translation-quotient}

Suppose a square-free edge set is invariant under translations by a
binary linear subspace \(H\) of index \(N=2^r=O(n)\). Then its cube-edge
density is \textbf{\(o(1)\)}, even if individual rules have high degree.

Write \(G=V/H\) and \(a_i=e_i+H\). The indicator for direction \(i\)
descends to \(G\) and is invariant under \(+a_i\). Within every class of
directions with the same \(a_i=a\), the active sets are pairwise
disjoint: simultaneous activity at one quotient point would give a
lifted square since both directions are also active at the opposite
quotient point. This holds also for \(a=0\). Thus
\(t_a=\sum_{i:a_i=a}p_i\le1\).

For each nonzero represented \(a\), pick the direction of largest
density; let \(m_a\) be the class size. Put \((x,x+a)\) in
\(D\subset G^2\) when this representative is active at \(x\). This is a
symmetric set disjoint from the diagonal. Its density is at least \[
N^{-1}\sum_{a\ne0:\,m_a>0}t_a/m_a\ \ge\ (np-1)_+^2/(nN).
\] For the second inequality use \(t_0\le1\) and Cauchy--Schwarz:
\((\sum_{a\ne0}t_a)^2\le n\sum_{a\ne0}t_a^2/m_a\le n\sum_{a\ne0} t_a/m_a\),
omitting empty classes. If \(p\) were bounded below on a subsequence,
then \(D\) would have positive density bounded away from zero, and
\(N\ge np\) would give \(r\to\infty\).

Identify \(G^2\) with length-\(r\) words on the alphabet
\(\{00,01,10,11\}\). The density Hales--Jewett theorem supplies a
combinatorial line in \(D\). Its four pairs are
\((x,y),(x+t,y),(x,y+t),(x+t,y+t)\) with \(t\ne0\) the indicator of the
variable coordinates. There are no loops, so \(a=x+y\) and \(b=x+y+t\)
are nonzero and distinct. The four edges alternate the chosen directions
for \(a,b\). Lifting from any preimage of \(x\) gives a forbidden
coordinate square. This proves the claim.

\subsection{Where the proof breaks for
cubics}\label{where-the-proof-breaks-for-cubics}

The two-color implication itself fails already on \(Q_5\) for cubics. In
coordinates \((x_1,x_2,s,v,w)\), set \(h=(1+v)(1+w)\) over
\(\mathbb F_2\) and retain only the first two directions, with \[
f_1=1+s+s x_2+(1+s)h,\qquad f_2=s+(1+s)x_1+h x_1.
\] When \(h=0\), the pair of functions is \((1,x_1)\) if \(s=0\) and
\((x_2,1)\) if \(s=1\). When \(h=1\), it is \((0,0)\) or \((x_2,1+x_1)\)
respectively. Thus no retained square exists, whereas the face at
\(s=v=w=0\) is entirely deleted. Both direction densities equal \(5/8\).
This does not disprove the density conjecture for cubic rules; it
pinpoints the failed implication.

\subsection{Verification and scope}\label{verification-and-scope}

Run the included Python 3 standard-library script, \textbf{verify.py}.
Truth-table bits are indexed by \(x=0,\dots,2^n-1\) in binary coordinate
order. In dimension four, all 2,048 possible quadratics were checked.
All 5,046 pairs of above-half own-direction-invariant functions were
tested for the six direction pairs; the 576 compatible pairs obey the
two-color conclusion. The four quadratic truth tables \textbf{cff3,
f5af, bbdd, 7e7e} (hexadecimal) define a graph with 24 of 32 edges and
\(b_i=1/2\) for all four directions. A direct scan confirms
square-freeness and equality in the main theorem. The script also checks
the tree optimum by independent dynamic programming and the
dimension-\(16\) iterated sign construction. Enumeration is a check; the
proof gives the arbitrary-dimensional statement.

\begin{itemize}
\tightlist
\item
  \href{https://www.erdosproblems.com/86}{Erdős problem \#86} and
  \href{https://www.erdosproblems.com/latex/86}{original references}.
\item
  Brass, Harborth and Nienborg,
  \href{https://doi.org/10.1002/jgt.3190190104}{On the maximum number of
  edges in a \(C_4\)-free subgraph of \(Q_n\)}, unrestricted lower
  constructions of excess order \(n^{-1/2}\). Baber,
  \href{https://arxiv.org/abs/1201.3587}{Turán densities of hypercubes},
  asymptotic unrestricted upper bound 0.60318.
\item
  Bialostocki previously obtained the unweighted \(1/2+1/(2\sqrt n)\)
  estimate when \textbf{both} color classes are initially required to
  avoid squares; see the
  \href{https://combinatorialpress.com/ars/vol16a/}{1983 journal record}
  and
  \href{https://www.its.caltech.edu/~dconlon/HypercubeCycle.pdf}{Conlon's
  statement of the corollary}. The distinctive step here is that
  quadraticity plus the one-sided prohibition force precisely that
  two-color condition on the above-half directions. The Walsh and
  weighted second-moment arguments are not claimed as new machinery.
\item
  Wang, \href{https://arxiv.org/abs/2611.14906}{The frustration index of
  fully frustrated hypercubes}, Nov.~11, 2026, treats a distinct fully
  frustrated subclass. Zhou et al.,
  \href{https://arxiv.org/abs/2611.27284}{The four-cycle extremal number
  of the seven-cube}, Nov.~20, 2026, proves the unrestricted finite
  value 304 for \(n=7\), not the asymptotic conjecture.
\item
  Polymath, \href{https://arxiv.org/abs/0910.3926}{A new proof of the
  density Hales--Jewett theorem}, used only in the translation-quotient
  theorem.
\item
  A targeted search did not locate the quadratic-direction implication
  or exact restricted formula; no priority or publication claim is made.
  For cubic rules the derivative-zero locus need not be affine, and a
  positively biased cubic need not have nonnegative conditional bias on
  every hyperplane. Extending the argument even to all cubics needs a
  new idea.
\end{itemize}

\end{document}
